% Source: Supplement: flow reductions and matching objective distinctions.
% Review and validation: docs/validation.md.
\section{网络流常用建模}
\label{math:flow-models}
\subsection*{二分图最大匹配与点容量}
将左右两侧点分别编号，连 $S\to L_i$ 容量 $1$、合法匹配边 $L_i\to R_j$ 容量 $1$、$R_j\to T$ 容量 $1$。最大流即最大匹配数；用残量反向容量读取匹配边是否被选择。某点可以匹配多次时，改其对应的源边或汇边容量。

顶点通行容量可用拆点处理：把 $v$ 拆为 $v_{in},v_{out}$，中间边容量设为点容量，所有入边接到 $v_{in}$，所有出边从 $v_{out}$ 发出。源汇是否需要限制容量按题意决定。

\subsection*{最大权匹配与最大权最大匹配}
匹配边费用设为负权值，其他边费用为零。直接跑最小费用最大流，得到的是\textbf{先最大化匹配数，再在这些匹配中最大化权和}。如果题目只要求最大权、允许少匹配，负权边会使这两个目标不同。

只求最大权时，可补一条 $S\to T$、容量为左侧点数、费用为零的旁路，再要求发送左侧点数单位流。旁路表示不匹配，固定总流量下的最小费用便对应任意大小的最大权匹配，最终权和为费用取负。若要求恰好 $K$ 对，则不加旁路，指定上限 $K$ 并检查能否达到。

\begin{lstlisting}[style=mtycpp,numbers=none]
MinCostFlow g;
g.init(6);
int s=5,t=6;
g.add(s,1,1,0),g.add(s,2,1,0);
g.add(3,t,1,0),g.add(4,t,1,0);
g.add(1,3,1,-7),g.add(2,4,1,3);
g.add(s,t,2,0);
auto res=g.flow(s,t,2);
if(res.ok&&res.flow==2){__int128 weight=-res.cost;}
\end{lstlisting}

\subsection*{最小割：二选一代价与闭合子图}
把点位于源侧或汇侧理解为一种二选一决策。$S\to v$ 的容量在 $v$ 放到汇侧时付出，$v\to T$ 的容量在 $v$ 放到源侧时付出。添加边 $u\to v$、容量为足够大的数，可禁止“选 $u$ 但不选 $v$”，因为该配置会切断这条边。

最大权闭合子图：每点收益为 $w_v$，若选 $u$ 必须选 $v$，就连 $u\to v$ 的大容量边；正收益从源连入，负收益向汇连出，容量分别为收益及其绝对值。答案为全部正收益之和减去最小割，残量源侧为选中点。

这里的“大容量”必须\textbf{严格大于所有可能有限割代价的上界}，并在流量类型和算法上限范围内。不能随意用一个看似很大的常量；费用和容量也不能互相替代。
